% ██████████████████████████████████████████████████████████████████
%
%  CHAPTER 1 — INTRODUCTION
%  Thesis: "Explainable AI in Multi-Tenant ERPs:
%           Navigating Data Sparsity and the Cold-Start Problem"
%
% ██████████████████████████████████████████████████████████████████
\chapter{Introduction}
\label{ch:introduction}

\section{Background}
\label{sec:background}

Over the past two decades, Enterprise Resource Planning (ERP) systems have evolved from monolithic, on-premise platforms reserved for multinational corporations into agile, cloud-native services that increasingly target Small and Medium-sized Enterprises (SMEs) \citep{haddara2011erp, gartner2023erp}. This democratization has been propelled chiefly by the rise of the multi-tenant Software-as-a-Service (SaaS) delivery model, in which a single hardened codebase serves hundreds or thousands of organizationally isolated customers from a shared infrastructure \citep{bezemer2010multi, krebs2014multitenancy}. The economic argument is compelling: SMEs gain access to enterprise-grade finance, inventory, customer relationship, and human resource modules at a fraction of the historical total cost of ownership, while vendors achieve economies of scale through consolidated operations.

In parallel with this architectural shift, the integration of Machine Learning (ML) into ERP suites has transitioned from an experimental differentiator to a baseline expectation \citep{idc2024ai}. Modern ERP platforms now ship with embedded predictive services such as late-payment risk scoring, sales-lead prioritization, employee-attrition forecasting, demand prediction, and expense-anomaly detection. These services are typically implemented using tabular supervised learning algorithms---most prominently Gradient Boosting machines, Random Forests, and regularized linear models---whose strong empirical performance on structured business data is well documented \citep{chen2016xgboost, lundberg2020local}. By translating raw transactional records into forward-looking probability estimates, these models promise to convert the ERP from a system of record into a system of insight.

This promise, however, is accompanied by a long-standing concern: the opacity of complex ML models. Ensemble methods such as Gradient Boosted Decision Trees aggregate hundreds of weak learners, producing predictions whose internal reasoning is not directly intelligible to domain experts, regulators, or end users \citep{rudin2019stop, lipton2018mythos}. In high-stakes business contexts---denying credit terms to a customer, deprioritizing a sales lead, or flagging an employee as a flight risk---opaque predictions are not merely inconvenient; they are operationally and ethically untenable. The field of Explainable Artificial Intelligence (XAI) has emerged to address precisely this gap, with two post-hoc, model-agnostic techniques rising to particular prominence: \textbf{SHAP} (SHapley Additive exPlanations), grounded in cooperative game theory \citep{lundberg2017unified}, and \textbf{LIME} (Local Interpretable Model-agnostic Explanations), which approximates the local decision surface with sparse linear surrogates \citep{ribeiro2016why}. Together, SHAP and LIME have become the de facto standard for justifying black-box predictions in commercial AI systems, and their adoption within ERP suites is increasingly viewed as a regulatory and trust prerequisite \citep{eu2024aiact}.

Despite this progress, a critical gap remains largely unaddressed in the literature on enterprise XAI. Published evaluations of SHAP and LIME overwhelmingly assume the \emph{data-rich regime} that characterizes large enterprises: hundreds of thousands of transactions, years of history, and a stable feature distribution. The multi-tenant SaaS reality is fundamentally different. In a typical mid-market ERP deployment, the long tail of tenants consists of newly onboarded SMEs whose databases contain only a few dozen to a few hundred invoices, leads, or employee records. For these tenants the predictive model itself is degenerate---incapable of separating signal from noise---yet the surrounding XAI layer continues to operate unchanged. SHAP will dutifully decompose a near-random prediction into precisely additive feature contributions, and LIME will fit a tidy linear surrogate around a decision boundary that has no statistical justification to exist. The user is presented with an authoritative-looking explanation built on top of what is empirically indistinguishable from a coin flip.

This research designates the phenomenon described above as the \textbf{Hallucinated Explanation} problem of multi-tenant XAI. It is a particularly pernicious failure mode because it fails silently and visibly at the same time: silently in the statistical sense, in that no error is raised and no warning is surfaced; visibly in the user-experience sense, in that the explanation is rendered with the same polished interface used for explanations derived from well-trained models. The cold-start tenant therefore receives the most dangerous combination possible---high apparent confidence supported by high apparent transparency---at precisely the moment when both are least warranted. To quantify this gap and motivate an architectural response, the present work conducts a controlled sparsity experiment across five simulated tenants ranging from 50{,}000 down to 50 training rows, and proposes a \textbf{Maturity-Aware AI Routing} pattern implemented inside a Django-based multi-tenant ERP.

% ── FIGURE 1.1 placeholder ────────────────────────────────────────
\begin{figure}[H]
  \centering
  % \includegraphics[width=0.85\textwidth]{figures/reliability_cliff_concept.pdf}
  \fbox{\parbox{0.85\textwidth}{\centering\vspace{2cm}
    \textit{[Placeholder for Figure~\ref{fig:reliability_cliff}]}\\[0.5em]
    Conceptual illustration of the ``Reliability Cliff''---SHAP explanation
    fidelity (Spearman rank correlation) plotted against tenant data volume
    on a logarithmic axis, contrasting the stable plateau exhibited by
    data-rich tenants with the sharp collapse observed in cold-start
    tenants below approximately one thousand training records.
    \vspace{2cm}}}
  \caption[Conceptual illustration of the Reliability Cliff in multi-tenant XAI.]%
  {Conceptual illustration of the \emph{Reliability Cliff} --- SHAP
   explanation fidelity (Spearman rank correlation against a data-rich
   reference tenant) degrades as tenant data volume drops along a
   logarithmic axis, contrasting data-rich tenants with cold-start
   tenants.}
  \label{fig:reliability_cliff}
\end{figure}

\section{Problem Statement}
\label{sec:problem_statement}

The trajectory described above gives rise to three tightly interconnected problems that together motivate this research:

\begin{enumerate}[label=\arabic*.]
  \item \textbf{Cold-Start Data Sparsity in Multi-Tenant ERPs.}
  The multi-tenant SaaS distribution model guarantees a permanent population of low-data tenants. New customers onboard continuously, and many SMEs operate at a transactional volume that never approaches the scale assumed by mainstream ML literature. Within such tenants, supervised learning algorithms cannot reliably estimate the underlying data-generating process: predictive performance, measured for instance by the Area Under the Receiver Operating Characteristic Curve (AUC-ROC), collapses toward the non-informative floor of $0.50$. This degradation has been observed empirically in the present work, where AUC decreased monotonically from $0.741$ at $50{,}000$ training rows to $0.500$ at $50$ rows. From a model-quality perspective, the predictions produced for such tenants carry no actionable information.

  \item \textbf{The ``Hallucinated Explanation'' Phenomenon.}
  SHAP and LIME are designed to \emph{explain} whichever model they are attached to; they do not, by construction, verify that the model itself is worth explaining. Consequently, both techniques continue to emit confident, well-formatted feature attributions even when the underlying classifier is statistically equivalent to random guessing. The empirical findings of this study quantify this failure mode for the first time in a multi-tenant ERP setting. The Spearman rank correlation between a sparse tenant's global SHAP ranking and a data-rich reference ranking degrades from $\rho = 1.000$ (50{,}000 rows) to $\rho = 0.917$ (5{,}000 rows), then crosses below the widely accepted $\rho = 0.7$ reliability threshold at the $1{,}000$-row mark ($\rho = 0.688$), and ultimately collapses to $\rho = 0.550$ at $50$ rows. Simultaneously, the variance of the absolute SHAP values---a direct proxy for explanation stability---spikes from $0.066$ in the data-rich regime to $0.176$ at $200$ rows, evidencing that sparse models \emph{grasp at straws}, concentrating their attribution mass on one or two spuriously dominant features. The result is an explanation layer whose surface fidelity is indistinguishable from that of a well-trained model, but whose epistemic content has been almost entirely eroded.

  \item \textbf{Absence of Architectural Safeguards.}
  Contemporary ERP and ML platforms expose XAI outputs through homogeneous user interfaces that draw no distinction between explanations grounded in a well-trained model and explanations generated atop a degenerate one. There is no widely adopted architectural pattern that conditions the \emph{delivery} of an explanation on the \emph{maturity} of the tenant producing it. As a consequence, the burden of validating explanation reliability is transferred implicitly to the end user---typically an SME finance officer or sales manager without the statistical training to perform that validation. The literature on responsible AI in enterprise systems emphasizes governance over algorithm choice and dataset documentation \citep{mitchell2019model, gebru2021datasheets}, but offers limited guidance on the systems-level controls required to prevent the delivery of unfaithful explanations in the first place.
\end{enumerate}

\noindent
These three problems collectively give rise to the following research questions, which this thesis seeks to answer:

\begin{enumerate}[label=\textbf{RQ\arabic*:}]
  \item To what extent does data sparsity quantitatively degrade the fidelity and stability of post-hoc XAI methods---specifically SHAP and LIME---in a multi-tenant ERP setting?
  \item Does a discernible threshold exist below which XAI explanations cease to be statistically reliable, and if so, where is that threshold located along the data-volume axis?
  \item How can the architecture of a multi-tenant ERP be designed such that the delivery of AI insights is automatically routed according to each tenant's data maturity?
  \item To what extent can a cold-start fallback mechanism, derived from cross-tenant global heuristics, preserve user trust and business utility when local explanations are suppressed?
\end{enumerate}

\section{Problem Limitation}
\label{sec:problem_limitation}

To preserve methodological rigor and the reproducibility of results, this research is intentionally bounded along four dimensions:

\begin{enumerate}[label=\arabic*.]
  \item \textbf{XAI Scope.} The investigation is restricted to two post-hoc, model-agnostic explanation families: SHAP---specifically the \texttt{TreeExplainer} for tree-ensemble models and the \texttt{LinearExplainer} for linear baselines---and LIME---specifically the \texttt{LimeTabularExplainer} for tabular feature spaces. Intrinsically interpretable models (e.g., decision rules, Generalized Additive Models), counterfactual explanations, and attention-based saliency techniques fall outside the scope of this work.

  \item \textbf{ML Model Scope.} The empirical evaluation targets three tabular supervised classification tasks representative of mainstream ERP analytics: Late-Payment Risk Prediction (finance module), Lead Scoring (CRM module), and Employee Attrition Prediction (HR module). The underlying algorithms are Gradient Boosting Classifiers and Random Forest Classifiers. Time-series forecasting, computer vision, and natural-language tasks are excluded.

  \item \textbf{Architectural Scope.} The multi-tenant architecture under study is implemented in Django~4.2 atop a PostgreSQL backend using the \emph{database-per-tenant} isolation strategy. Alternative isolation strategies---shared-schema row-level security or schema-per-tenant---and alternative web frameworks fall outside the scope.

  \item \textbf{Experimental Scope.} Tenant sparsity is simulated by stratified data slicing of open-source datasets---the Brazilian Olist e-commerce corpus for the Late-Payment task and the IBM HR Employee Attrition dataset for the Attrition task---rather than by longitudinal observation of live SaaS customers. The simulated tenants therefore approximate cross-sectional data-volume regimes but do not capture potential concept drift introduced by the temporal onboarding trajectories of real tenants.
\end{enumerate}

\section{Research Objectives}
\label{sec:objectives}

In response to the research questions formulated above, this thesis pursues five interconnected objectives:

\begin{enumerate}[label=\arabic*.]
  \item \textbf{Quantify XAI degradation under data sparsity} by training the same supervised classifier on five tenant slices of geometrically decreasing volume ($50{,}000$; $5{,}000$; $1{,}000$; $200$; $50$ rows) and measuring the corresponding evolution of explanation fidelity (Spearman rank correlation), local explanation stability (LIME top-$k$ Jaccard agreement), global explanation concentration (Shannon entropy of normalized SHAP importances), and explanation volatility (variance of absolute SHAP values).

  \item \textbf{Identify a practical cold-start threshold} expressed in terms of the minimum number of historical records required for SHAP-derived global rankings to remain within an a-priori reliability band ($\rho \geq 0.7$) relative to a data-rich reference tenant.

  \item \textbf{Design a ``Maturity-Aware AI Routing'' architectural pattern} that conditions the delivery of XAI outputs on tenant maturity, formally separating two service tiers: a local-prediction tier for tenants whose data volume exceeds the empirically derived threshold, and a global-heuristic tier for tenants below it.

  \item \textbf{Implement and evaluate the tier-logic gate} as a first-class component of the Django-based ERP service layer (\texttt{LatePaymentPredictor.predict\_or\_heuristic}) and demonstrate that the gating decision is auditable, side-effect-free, and integrable with existing prediction pipelines without breaking backward compatibility.

  \item \textbf{Synthesize cross-tenant global heuristics} from the data-rich tenant's SHAP attributions and persist them as a versioned JSON artifact (\texttt{cold\_start\_heuristic.json}), thereby providing cold-start tenants with platform-wide business guidance in lieu of unreliable local explanations.
\end{enumerate}

\section{Research Benefits}
\label{sec:benefits}

The contributions of this research are organized along two complementary axes, in the manner conventional to applied informatics scholarship.

\subsection{Theoretical Benefits}
\label{subsec:theoretical_benefits}

\begin{itemize}
  \item The study provides, to the best of the author's knowledge, the \textbf{first empirical characterization} of SHAP and LIME degradation profiles in a controlled multi-tenant ERP setting, expressed through quantitative metrics---Spearman fidelity, top-$k$ Jaccard stability, and SHAP attribution variance---rather than the qualitative case studies that dominate the existing literature.

  \item It contributes the notion of an \textbf{explanation reliability cliff}---an empirically grounded data-volume threshold below which post-hoc XAI methods cease to be epistemically trustworthy---and locates that threshold within the regime relevant to the multi-tenant SaaS economy.

  \item It reframes the academic position of post-hoc XAI methods, advocating their adoption not merely as a user-facing reporting feature but as a \textbf{diagnostic validation instrument}: a test of whether a deployed model deserves to be deployed at all.

  \item It introduces a vocabulary---\emph{Hallucinated Explanation}, \emph{Maturity-Aware AI Routing}, \emph{Cold-Start Heuristic Tier}---that future work on responsible AI in low-data B2B environments may build upon.
\end{itemize}

\subsection{Practical Benefits}
\label{subsec:practical_benefits}

\begin{itemize}
  \item For \textbf{software architects and platform engineers}, the work delivers a directly reusable architectural pattern, complete with an executable reference implementation in Django, that can be integrated into existing multi-tenant ML pipelines with minimal disruption to upstream callers.

  \item For \textbf{SME end users} of multi-tenant ERP systems, the proposed safeguards reduce the risk of basing material business decisions---extending or denying credit, prioritizing or deprioritizing sales leads, intervening on suspected attrition risks---on AI explanations that have no statistical foundation.

  \item For \textbf{SaaS vendors}, the cold-start heuristic tier offers a scalable mechanism to extend the perceived value of AI features to newly onboarded customers without exposing the vendor to the reputational and regulatory risk associated with serving hallucinated explanations.

  \item For \textbf{regulators and auditors} operating in jurisdictions governed by emerging AI-accountability frameworks (e.g., the EU AI Act \citep{eu2024aiact}), the architecture provides a concrete, auditable control point at which the suppression of unreliable explanations is enforced by the system itself, rather than relegated to user discretion.
\end{itemize}

\noindent
Taken together, these theoretical and practical contributions position the present thesis as both a diagnostic study of an under-examined failure mode and a constructive blueprint for mitigating that failure mode at the architectural level of multi-tenant enterprise software.
